\section{模型的评价与推广}

\subsection{模型优点}

\textbf{（1）用一条恒等式统一了四问的经济逻辑。}式~\eqref{eq:identity} 把执行缺额写成预测误差减去计划弃光量，说明在“严格按计划执行”的口径下，紧急购电只由预测精度和保险水平决定，与储能调度无关。这条恒等式使问题二的“买多少保险”和问题三的“预报值多少钱”变成两个可以独立处理、且都能用闭式公式刻画的子问题。

\textbf{（2）安全裕度有理论最优解而非经验调参。}紧急购电价的 5 倍非对称代价使最优保险量恰好等于预测误差分布的 $0.8$ 分位点。实测全年费用曲线的最低点落在 $q=0.80$，且在 $m=2,3,8,10$ 倍时 $q=1-1/m$ 始终最优，理论预测与数值实验完全吻合，模型的结论不是对某一组数据的过拟合。

\textbf{（3）严格因果且可复现。}问题二至四的每一天都只使用决策时刻之前的信息，全流程由 \texttt{run\_all.py} 一键复现（约 30 秒），所有数值都可由结果文件回读核对。

\subsection{模型缺点}

\textbf{（1）储能严格按计划充放电，牺牲了执行阶段的自适应能力。}这也是式~\eqref{eq:identity} 成立的前提。若允许储能根据实际缺额临时放电，紧急购电可以进一步下降，代价是充放电量偏离计划、SOC 轨迹不再受控。当前口径下电池的应急缓冲作用只能通过“多买”来替代，这是保守的处理。

\textbf{（2）日循环约束放弃了跨日储能。}由于电价日型固定，跨日套利空间本就很小（灵敏度分析显示差值在 $10^{-3}$ 量级），但在连续阴天等场景下放开该约束仍有收益。

\textbf{（3）附件 3 的预报没有被充分利用。}表 7 显示附件 3 的预报误差约为历史均值的 $2.4$ 倍，若把两路信息按历史表现加权融合，预测误差还能进一步下降。本文为了严格对应题面给定的信息结构，没有做这一步。

\subsection{推广方向}

本文的“预测误差 + 报童裕度 + 滚动修正”框架不限于微网购电。凡是满足以下三条的场景都可以直接套用：存在一个按计划结算的基础量、偏差按非对称价格事后结算、并且有可逐步更新的预测。典型的例子包括工业用户的需量申报与偏差考核、电力市场的日前申报与实时平衡、以及供热的以热定电与调峰补偿。模型中唯一需要替换的部分是预测模块——把 7 日历史均值替换为该场景的专用预测即可，安全裕度的分位点仍由 $q=1-1/m$ 给出，其中 $m$ 是偏差惩罚倍数。
